\scanpage{244}\mspage{- 15 -}
n$^{\mathrm{o}}$ 227

linéaire, le système des hypersurfaces de degré donn\handwritten{\'}{e} sur une
variété projective.

Soit donc $V$ une variété projective, et $R=k[x]=k[x_{\unsure{0|o}},\ldots,x_N]$
son anneau de coordonnées homogènes, comme dans IAG, ch.IV. Soit $I$
\deleted{la fermeture intégrale de $R$}, et $R_n$, $I_n$ les éléments homogènes de
degré $n$ de $R$ et de $I$. Soit $f\handwritten{\in}R_n$. On va associer un diviseur positif
$Z(f)$ à $f$, qui sera appelé l'hypersurface section de $V$ par $f$. Soit
$W^{r-\handwritten{1}}$ un diviseur premier de $V$. Supposons que $W$ ne soit pas contenu
dans l'hyperplan $X_0=0$. Alors $f(x)/x_0^n$ est une fonction $\handwritten{\varphi}$ sur $V$. On
définit la multiplicité de $W$ dans le diviseur $Z(f)$ comme étant
$\operatorname{ord}_W\handwritten{\varphi}$. (Il est clair que $\operatorname{ord}_W\handwritten{\varphi}$ ne dépend pas de l'hyperplan auxiliaire,
et qu'on obtient bien un diviseur positif). Les assertions
suivantes découlent trivialement des définitions :

Tous les cycles $Z(f)$, $f\handwritten{\in}R_n$, sont linéairement équivalent\unsure{:|.} En
fait, si $f$, $g\handwritten{\in}R_n$, on a $(f/g)=Z(f)-Z(g)$.

Si $f_0\handwritten{\in}R_m$, alors l'ensemble des fonctions $f/f_0$, $f\handwritten{\in}R_m$ est
un module $L_0$, et le sup des diviseurs polaires de ces fonctions est
exactement $Z(f_0)$. L'ensemble des diviseurs $Z(f)$ où $f$ parcourt $R_n$ est
un système linéaire, définit par le module $L_0$. On voit aussi qu'il
n'a pas de composantes fixes.

Le système linéaire complet dérivé du système d'hypersurface
de degré $n$ est caractérisé algébriquement par le théorème suivant
(Zariski, th.1 de ES, démonstration facile et élémentaire ) :

\margin{\handwritten{?}}TH\handwritten{\'}{E}OR\handwritten{\`}{E}ME 2. - \uline{Si} $\handwritten{\varphi}$ \uline{est une fonction telle que} $(\handwritten{\varphi})\handwritten{\geq}-Z(f)$ \uline{pour}
$f\handwritten{\in}R_n$, \uline{alors il existe} $g\handwritten{\in}I_n$ \uline{tel que} $\handwritten{\varphi}=g/f$. \uline{Si} $V$ \uline{est normale},
\uline{alors le système d'hypersurfaces de degré} $n$ \uline{est complet pour tout} $n$
\uline{suffisam\inserted{m}ent grand}. (Voir la prop.11 de IAG, ch.IV).
% archive sheet number: 243
